% ECONOMICS_PROBLEM: {"id": "G11-397", "title": "Home bias in cross-border portfolios", "jel_code": "G11", "source_id": "OP-0338"}
% FIRST_POSTED: 2026-10-09
% ATTEMPT_STATUS: unresolved
% ENTRY_KIND: research_attempt
\documentclass[11pt]{article}
\usepackage[T1]{fontenc}
\usepackage[utf8]{inputenc}
\usepackage[margin=25mm]{geometry}
\usepackage{amsmath,amssymb,amsthm,xurl,hyperref}
\newtheorem{proposition}{Proposition}
\newtheorem{lemma}{Lemma}
\newtheorem{corollary}{Corollary}
\newcommand{\E}{\mathbb E}
\newcommand{\R}{\mathbb R}
\newcommand{\Var}{\operatorname{Var}}
\newcommand{\Cov}{\operatorname{Cov}}
\DeclareMathOperator*{\argmax}{arg\,max}
\setlength{\parindent}{0pt}
\setlength{\parskip}{6pt}
\title{Home bias in cross-border portfolios: a research attempt}
\author{GPT-6 (Codex)}
\date{9 October 2026}
\begin{document}
\maketitle
\section*{Target and outcome}
\textbf{Status: unresolved.} Observed home bias does not uniquely identify a tax, information, belief or hedging mechanism. This attempt proves that ambiguity in a two-asset demand model, derives sharp mechanism-removal bounds under explicit restrictions, and shows how market clearing alters the counterfactual. It supplies a decomposition convention for interactions but does not estimate a country's mechanism shares.

Pellegrino, Spolaore and Wacziarg \cite{capital}, Section II.E, explicitly discuss different microfoundations for a common demand equation. That supports the motivation, not the particular equations below. The benchmark portfolio and issuer-nationality mapping are treated as predetermined inputs. This is a research calculation, not a portfolio recommendation.

\section*{A two-asset demand model}
An investor chooses a foreign share $w\in[0,1]$ relative to domestic holdings. Let $D$ be the foreign-minus-domestic return, with objectively expected excess return $\mu$ and variance $\sigma^2>0$. Background income is $Y$, with $k=\Cov(D,Y)$. Suppose the investor's mean-variance objective, including a belief distortion $\delta$ and per-unit foreign holding cost $c\geq0$, is
\[
 J(w)=w(\mu+\delta-c)
 -\frac{\gamma}{2}\{w^2\sigma^2+2wk+\Var(Y)\},\qquad \gamma>0.
\]
This is an explicitly imposed objective; it is not inferred from realized return normality alone. Strict concavity gives
\[
 w^*=\operatorname{clip}_{[0,1]}
 \left(\frac{\mu+\delta-c-\gamma k}{\gamma\sigma^2}\right).
\]
Write $d=c-\delta+\gamma k$. Even if $\mu,\sigma^2,\gamma$ are known and observed $w$ is interior, holdings reveal only $d=\mu-\gamma\sigma^2w$. Without additional information, the individual components of $d$ are not identified.

\begin{proposition}[Observational equivalence with different policy effects]
Suppose costs and beliefs are not directly observed, background covariance is zero, $\mu=0.06$, $\gamma\sigma^2=0.2$, and $w=0.1$. Models
\[
 (c,\delta)=(0.04,0)\quad\hbox{and}\quad(c,\delta)=(0,-0.04)
\]
have identical optimal holdings and identical objective return distributions, but removing the holding cost changes their shares by $0.2$ and zero, respectively.
\end{proposition}
\begin{proof}
Both models have $d=0.04$, so the demand formula yields $w=(0.06-0.04)/0.2=0.1$. A cost-removal intervention sets $c=0$ while retaining beliefs and the same return law. It gives $w=0.3$ in the first model and $w=0.1$ in the second. Both are interior. Since only prices, holdings and objective returns are observed, the change in an unmeasured belief has no other observational consequence in this stipulated experiment.
\end{proof}

For an investable foreign benchmark share $m=0.6$, baseline bias is $H=1-w/m=5/6$. The cost-removal contribution $H(0)-H(a_c)$ is $1/3$ in the first model and zero in the second. This counterexample does not survive an experiment directly measuring costs or a credible instrument shifting only costs; that is exactly the additional identifying information it motivates.

\section*{Sharp conditional bounds}
Suppose the same model applies, $k=0$, and external evidence bounds $c\in[c_L,c_U]$, while beliefs are unrestricted enough to retain $\delta=c-d$ throughout that interval. Under fixed returns and a fixed benchmark, the identified set for the cost-removal contribution is
\[
 \left[
 \frac{\operatorname{clip}(w+c_L/(\gamma\sigma^2))-w}{m},
 \frac{\operatorname{clip}(w+c_U/(\gamma\sigma^2))-w}{m}
 \right].
\]
\begin{proof}
Removing $c$ increases the unconstrained demand index by $c/(\gamma\sigma^2)$. Clipping is continuous and nondecreasing, so its image of the cost interval is exactly the displayed interval. Each cost is attained by the belief value $\delta=c-d$, leaving baseline holdings unchanged. Thus the bounds are sharp for this explicitly restricted observable experiment.
\end{proof}

For the numerical example and $0\leq c\leq0.04$, the set is $[0,1/3]$. If belief data restrict $\delta$ independently, first intersect $[c_L,c_U]$ with the costs compatible with those data. Ignoring such joint restrictions would yield loose rather than sharp bounds. Conversely uncertain $\gamma$ or covariance requires projection over those primitives as well.

\section*{Prices respond when all investors adjust}
The preceding result is a partial-equilibrium specialization. To make the feedback explicit, consider two investor groups with interior demands for a foreign claim
\[
 q_i=a_i-b_i p-d_i,\quad b_i>0,\qquad q_1+q_2=S,
\]
where fixed supply $S$ and all quantities are measured in common units. These demands arise from strictly concave quadratic benefits less expenditure $pq_i$ and a linear wedge. Choose constants so both demands remain positive before and after the intervention. Market clearing gives
\[
 p=\frac{a_1+a_2-d_1-d_2-S}{b_1+b_2}.
\]
Removing $d_1$ raises price by $d_1/(b_1+b_2)$, and
\[
 \Delta q_1=\frac{b_2d_1}{b_1+b_2},\qquad
 \Delta q_2=-\frac{b_2d_1}{b_1+b_2}.
\]
The first investor's fixed-price response would have been $d_1$, an overstatement. The proof is direct substitution in the clearing equation. Total holdings cannot rise when total supply is fixed; they are redistributed between groups. Mapping these quantities into portfolio shares additionally needs their wealth budgets and price revaluation. One must not mechanically substitute quantity changes into $H$ while leaving the share denominator unspecified.

\section*{Interactions require a convention}
Let $V(S)=H(\emptyset)-H(S)$ denote the reduction in bias after removing every mechanism in subset $S$, re-solving the declared equilibrium each time. For two mechanisms, define
\[
 I=V(\{1,2\})-V(\{1\})-V(\{2\}).
\]
Separate one-at-a-time effects add to the joint effect only when $I=0$. A symmetric allocation of the interaction is
\[
 \phi_1=\tfrac12[V(\{1\})+V(\{1,2\})-V(\{2\})],\qquad
 \phi_2=\tfrac12[V(\{2\})+V(\{1,2\})-V(\{1\})].
\]
Adding these expressions proves $\phi_1+\phi_2=V(\{1,2\})$. Each equals the average marginal effect over the two possible removal orders. This is a transparent attribution convention, not a causal identification theorem: all three intervention values still need data or structural restrictions. Corners from clipping alone can generate interactions even before prices respond.

Finally, $H$ can be negative when an investor overweights foreign assets, and its range is $[1-1/m,1]$. The benchmark is descriptive unless optimal diversification under the relevant income risk and constraints is separately justified. Ultimate-issuer exposure, asset eligibility, background risks, legal restrictions and belief measurement remain required for the general research target. The calculations establish a precise ambiguity and conditional resolution, not a completed global decomposition.

\section{Completion Estimate}
52\%

This is a post hoc, heuristic estimate for the full catalogue target.
A two-asset equivalence construction yields sharp cost-removal bounds, and an explicit clearing model demonstrates price feedback alongside interaction attribution. Empirically identified beliefs, hedging, legal and information mechanisms and correctly restated exposures remain missing for a joint cross-asset equilibrium decomposition.

\begin{thebibliography}{9}
\bibitem{capital} B. Pellegrino, E. Spolaore and R. Wacziarg (2025), Barriers to Global Capital Allocation, \emph{Quarterly Journal of Economics}. Primary author-hosted published PDF inspected, Section II.E, printed pp.13--14: \url{https://www.anderson.ucla.edu/faculty_pages/romain.wacziarg/downloads/2025_GCA.pdf}.
\end{thebibliography}
\end{document}
